\section*{अध्याय \ref{ch:b3:innate-immunity} --- जन्मजात प्रतिरक्षा और शोथ}

\begin{solution}{exo:b3:innate-immunity:1}
\omterm{def:b3:innate-immunity:innate}{अवरोध}: त्वचा का किरेटिन (यांत्रिक), वायुमार्ग की श्लेष्मा तथा \omterm{def:b3:cytoskeleton-motility:cilia}{पक्ष्माभ}
(फँसाना और सफ़ाई), आमाशय का अम्ल (निगले सूक्ष्मजीव मारता है), स्रावों
में लाइसोज़ाइम और डिफ़ेंसिन (लयन तथा छेदन), सहभोजी \omterm{def:b3:microbiomes:symbiosis}{माइक्रोबायोम}
(परिवेश-कोटर घेरता है)। कोशिकाएँ: \omterm{def:b3:innate-immunity:innate}{न्यूट्रोफिल} (\omterm{def:b3:innate-immunity:killing}{भक्षण} से जीवाणु मारते
हैं), \omterm{def:b3:innate-immunity:innate}{महाभक्षक} (खाते, सफ़ाई करते, ख़तरे की घंटी बजाते हैं), \omterm{def:b3:innate-immunity:innate}{वृक्षाभ
कोशिकाएँ} (प्रतिजन लसीका ग्रंथियों तक ले जाती हैं और अनुकूली अनुक्रियाएँ
शुरू करती हैं), मास्ट कोशिकाएँ (हिस्टामिन छोड़ती हैं, \omterm{def:b3:innate-immunity:inflammation}{शोथ} करती हैं),
\omterm{def:b3:innate-immunity:innate}{प्राकृतिक घातक कोशिकाएँ} (संक्रमित और रूपांतरित कोशिकाएँ मारती हैं)।
\end{solution}

\begin{solution}{exo:b3:innate-immunity:2}
सूक्ष्मजीवों के पूरे वर्ग में साझा, उनके लिए अनिवार्य (ताकि बचने के लिए
छोड़ा न जा सके), और परपोषी में अनुपस्थित। \omterm{def:b3:bacteriology:envelope}{लिपोपॉलिसैकेराइड} --- TLR4;
फ्लैजेलिन --- TLR5; द्विलड़ी RNA --- TLR3 (और कोशिकाद्रव्य में
RIG-I); अमेथिलित \omterm{def:b3:chromatin-epigenetics:methylation}{CpG} DNA --- TLR9; पेप्टाइडोग्लाइकन के टुकड़े --- NOD
प्रोटीन; $\beta$-ग्लूकन --- डेक्टिन-1।
\end{solution}

\begin{solution}{exo:b3:innate-immunity:3}
लाली: हिस्टामिन, नाइट्रिक ऑक्साइड और प्रोस्टाग्लैंडिनों से अणुधमनियों
का फैलाव। ऊष्मा: वही बढ़ा हुआ रक्त-प्रवाह। सूजन: अंतःस्तरी संकुचन के
बाद अणुशिराओं से प्लाज़्मा का रिसाव (हिस्टामिन, C3a, C5a)।
पीड़ा: ब्रैडीकाइनिन और प्रोस्टाग्लैंडिन E$_{2}$, जो पीड़ाग्राहियों को
संवेदी बना देते हैं, तथा सूजन का दबाव।
\end{solution}

\begin{solution}{exo:b3:innate-immunity:4}
चिरप्रतिष्ठित: किसी पृष्ठ पर प्रतिरक्षी (या \omterm{prop:b3:innate-immunity:systemic}{C-क्रियाशील प्रोटीन}) से
बँधता C1q। लेक्टिन: सूक्ष्मजीवी शर्कराओं पर मैनोस-बंधक लेक्टिन।
वैकल्पिक: स्वतःस्फूर्त C3 स्वतःचाल और किसी भी असुरक्षित पृष्ठ पर C3b का
जमना। परिणाम: C3b \omterm{def:b3:innate-immunity:killing}{भक्षण} के लिए \omterm{def:b3:innate-immunity:complement}{ऑप्सोनीकरण} करता है; C3a और C5a \omterm{def:b3:innate-immunity:inflammation}{शोथ} तथा
भर्ती करते हैं; C5b वह \omterm{def:b3:innate-immunity:complement}{झिल्ली-आक्रमण संकुल} शुरू करता है जो लयन करता है।
\end{solution}

\begin{solution}{exo:b3:innate-immunity:5}
$B = [k_{0}/(a-d)](e^{(a-d)t} - 1)$, जिसमें $a - d = 0.1$: \qty{60}{s} पर
$20(e^{6} - 1) \approx 8000$; \qty{90}{s} पर $20(e^{9} - 1) \approx
1.6\times 10^{5}$। परपोषी कोशिका: $k_{0}/(d - a) = 2/0.1 = 20$।
\end{solution}

\begin{solution}{exo:b3:innate-immunity:6}
\omterm{def:b3:innate-immunity:inflammation}{सिलेक्टिनों} पर लुढ़कना (कार्बोहाइड्रेट बाँधते लेक्टिन); कीमोकाइन-संकेतन
जो \omterm{def:b3:innate-immunity:inflammation}{इंटिग्रिन} सक्रिय करता है; अंतःस्तरी आसंजन अणुओं (ICAM) से
\omterm{def:b3:innate-immunity:inflammation}{इंटिग्रिनों} का दृढ़ आसंजन; अंतःस्तरी कोशिकाओं के बीच से पार जाना;
प्रवणता के सहारे रसोसंचलन। $\beta_{2}$ \omterm{def:b3:innate-immunity:inflammation}{इंटिग्रिनों} के बिना \omterm{def:b3:innate-immunity:innate}{न्यूट्रोफिल}
लुढ़कते तो हैं पर न कभी रुकते हैं न पार जाते: वे रक्त में जमा हो जाते
हैं (बहुत ऊँची गणना), संक्रमण बिना पूय बने लौटते रहते हैं, घाव खराब
भरते हैं, और नाभि-रज्जु देर से अलग होती है।
\end{solution}

\begin{solution}{exo:b3:innate-immunity:7}
चिरकारी कणिकागुल्मी रोग: \omterm{def:b3:innate-immunity:killing}{NADPH ऑक्सिडेज़} दोषपूर्ण है, इसलिए भक्षकाय में
न सुपरऑक्साइड बनता है, न हाइड्रोजन परॉक्साइड, न हाइपोक्लोराइट। उन जीवों
से संक्रमण जो अपना ही हाइड्रोजन परॉक्साइड नष्ट कर देते हैं
(कैटालेज़-धनात्मक): \emph{Staphylococcus aureus}, \emph{Aspergillus},
\emph{Burkholderia}, \emph{Serratia}, \emph{Nocardia}। कणिकागुल्म
इसलिए बनते हैं कि \omterm{def:b3:innate-immunity:innate}{महाभक्षक} ऐसे सूक्ष्मजीव निगल लेते हैं जिन्हें वे मार
नहीं सकते और उन्हें चुनकर घेर देते हैं, लगातार भर्ती के साथ, किसी
चिरकारी घाव में।
\end{solution}

\begin{solution}{exo:b3:innate-immunity:8}
निरोधी ग्राही MHC~I गिनते हैं; सक्रियकारी ग्राही दबाव-बंधक गिनते हैं;
कोशिका तब मरती है जब सक्रियण निरोध पर भारी पड़े। (क) MHC~I खोने से
निरोध हट जाता है: \omterm{def:b3:innate-immunity:innate}{प्राकृतिक घातक कोशिका} उसे मार देती है --- यही
T~कोशिकाओं से बचने की कीमत है। (ख) MHC~I रखे रहने पर भी, यदि दबाव-बंधक
प्रदर्शित हों तो सक्रियण फिर भी जीत सकता है और कोशिका मार दी जाती है;
इसीलिए \omterm{def:b3:cancer-biology:hallmarks}{अर्बुद} प्रायः दबाव-बंधक झाड़ देते या दबा देते हैं।
\end{solution}

\begin{solution}{exo:b3:innate-immunity:9}
शुद्ध प्रोटीन कोई रोगजनक-नमूना नहीं ढोता, इसलिए जो \omterm{def:b3:innate-immunity:innate}{वृक्षाभ कोशिकाएँ} उसे
उठाती हैं वे परिपक्व नहीं होतीं, उसे बिना सह-उद्दीपन के प्रस्तुत करती
हैं, और जो T~कोशिकाएँ उसे पहचानती हैं वे बंद कर दी जाती हैं या उसे
अनदेखा कर देती हैं। कोई सहवर्धक --- फिटकरी, कोई लिपिड, कोई \omterm{def:b3:innate-immunity:prr}{टोल-सदृश
ग्राही} बंधक --- नमूना-ग्राही छेड़ता है, \omterm{def:b3:innate-immunity:innate}{वृक्षाभ कोशिका} परिपक्व करता है
और दूसरा संकेत जुटाता है। जेनवे की बात यही है: अनुकूली तंत्र किसी
प्रतिजन पर अनुक्रिया तभी देता है जब जन्मजात तंत्र कहे कि वह प्रतिजन
किसी सूक्ष्मजीव के साथ आया था।
\end{solution}

\begin{solution}{exo:b3:innate-immunity:10}
ज्वर के बिना $48$ द्विगुणन: $10^{4}\times 2^{48} \approx 3\times
10^{18}$ (कोई ऐसी बेतुकी संख्या जो दिखाती है कि परिणाम वृद्धि नहीं,
मारना तय करता है)। ज्वर के साथ $32$ द्विगुणन: $4\times 10^{13}$ ---
यानी \omterm{def:b3:innate-immunity:innate}{न्यूट्रोफिलों} के सामने $2^{16} \approx 65\,000$ गुना कम जीवाणु।
लागत: दिन में $3\times 0.12\times 8 = \qty{2.9}{MJ}$, यानी कोई बड़ा
भोजन। तब लायक, जब रोगजनक धीमा पड़े और परपोषी के पास भंडार हों; तब नहीं,
जब रोगजनक \qty{40}{\celsius} पर भी उतना ही बढ़े, या परपोषी भूखा हो, या
ज्वर स्वयं ख़तरनाक तापमानों के पास पहुँचने लगे।
\end{solution}

\begin{solution}{exo:b3:innate-immunity:11}
कारक~H बाँधना परपोषी का अपना नियामक भर्ती कर लेना है: $d$ चढ़कर $a$ से
ऊपर हो जाता है और वह पाश उस सूक्ष्मजीव पर वैसे ही क्षीण होता है जैसे
किसी परपोषी कोशिका पर। C3b काटना जमी कन्वर्टेज़ हटा देता है: फिर से
$d$ में बढ़त। कोई संपुट ऐसा पृष्ठ देता है जिस पर कन्वर्टेज़ ठीक से बनती
ही नहीं, यानी $a$ घटता है, और जो C3b जमता भी है उसे भक्षककोशिका के
ग्राहियों से छिपा देता है। सबसे सस्ता: कोई एक पृष्ठीय प्रोटीन जो
कारक~H बाँध ले, जिसे परपोषी मुफ़्त जुटाता है और जो काम कर देता है।
\end{solution}

\begin{solution}{exo:b3:innate-immunity:12}
किसी किरच पर वह फैलाव, रिसाव और भर्ती कुछ मिलीलीटर वाहिकाओं को शामिल
करते हैं और बचाव वहीं पहुँचाते हैं जहाँ उसकी ज़रूरत है; वही मध्यस्थ पूरे
परिसंचरण में छूट जाएँ तो हर वाहिका फैलती है और हर केशिका रिसती है,
इसलिए दबाव ढह जाता है और हर जगह स्कंदन छिड़ जाता है --- शरीरक्रिया वही
है, पैमाना घातक। TNF आमवाती संधिशोथ का संधिकला-शोथ चलाता है, इसलिए उसे
रोकना रोग से राहत देता है; पर TNF वही है जो \omterm{def:b3:innate-immunity:innate}{महाभक्षकों} को उन
कणिकागुल्मों में संगठित रखता है जो प्रसुप्त तपेदिक-दंडाणु घेरे रहते
हैं, और उसे रोकना प्रसुप्त तपेदिक को छूट जाने देता है।
\end{solution}

\begin{solution}{pb:b3:innate-immunity:1}
\textbf{1.} एक घंटा $1.5$ द्विगुणन है: $10^{4}\times 2^{1.5} \approx
2.8\times 10^{4}$।
\textbf{2.} $2\times 10^{4}$ \omterm{def:b3:innate-immunity:innate}{न्यूट्रोफिल} $4\times 10^{5}$ तक मारते हैं;
जीवाणु केवल $2.8\times 10^{4}$ से $8\times 10^{4}$ तक बढ़ते हैं:
मारना वृद्धि से पाँच गुना आगे है।
\textbf{3.} घंटा 2: मारने से पहले $8\times 10^{4}$, बाद में एक भी नहीं।
समष्टि \qty{2}{h} के आसपास लगभग $8\times 10^{4}$ पर शिखर छूती है और
उसके बाद शून्य है; \qty{6}{h} पर एक भी नहीं बचता।
\textbf{4.} \qty{3}{h} पर आगमन: तब $10^{4}\times 2^{4.5} = 2.3\times
10^{5}$। घंटा 4: $6.4\times 10^{5} - 4\times 10^{5} = 2.4\times
10^{5}$; घंटा 5: $6.8\times 10^{5} - 4\times 10^{5} = 2.8\times 10^{5}$;
घंटा 6: $7.9\times 10^{5} - 4\times 10^{5} = 3.9\times 10^{5}$ और
चढ़ता हुआ --- \omterm{def:b3:innate-immunity:innate}{न्यूट्रोफिल} अब साथ नहीं दे पाते, और कोई फोड़ा बन जाता है।
दो घंटे की देरी किसी इलाज को चिरकारी संक्रमण में बदल देती है।
\textbf{5.} $8\times 10^{4}/20 = 4000$ \omterm{def:b3:innate-immunity:innate}{न्यूट्रोफिल} मारते हुए मरते हैं,
साथ ही वे $2\times 10^{4}$ जो पहुँचे थे और दिन भर में वैसे भी मर जाते
हैं। पूय मरे \omterm{def:b3:innate-immunity:innate}{न्यूट्रोफिल}, मरे जीवाणु और द्रवित ऊतक है; \omterm{def:b3:innate-immunity:innate}{महाभक्षक} उन
लाशों को खाकर उसे साफ़ करते हैं, वरना उसे बहना पड़ता है।
\textbf{6.} प्रति घंटे $2000$ आगमन, जो $4\times 10^{4}$ मारते हैं:
घंटा 2, $8\times 10^{4} - 4\times 10^{4} = 4\times 10^{4}$; घंटा 3, $1.1\times
10^{5} - 4\times 10^{4} = 7\times 10^{4}$; घंटा 4, $2\times 10^{5} -
4\times 10^{4} = 1.6\times 10^{5}$ --- बिना किसी
सीमा के बढ़ता हुआ। ऐसा रोगी किसी मामूली टीके को भी नहीं थाम सकता,
इसलिए पहले ही संकेत से मारने का काम \omterm{def:b3:bacteriology:antibiotics}{प्रतिजैविकों} को करना पड़ता है।
\textbf{7.} $B = 12.5(e^{0.08t} - 1)$: \qty{30}{s}, $125$; \qty{60}{s},
$1500$; \qty{120}{s}, $1.8\times 10^{5}$।
\textbf{8.} $e^{0.08t} = 1.6\times 10^{5}$: $t = 12/0.08 = \qty{150}{s}$।
\textbf{9.} $k_{0}/(d - a) = 10$।
\textbf{10.} $B(120) = 33(e^{3.6} - 1) \approx 1200$, यानी $150$ गुना कम।
वह संपुट मिनट ख़रीदता है --- बढ़ने का समय, और तब तक की सुरक्षा जब तक
प्रतिरक्षी आकर उसी संपुट पर चिरप्रतिष्ठित पथ शुरू न कर दें।
\textbf{11.} $10^{4}\times 2\times 10^{6} = 2\times 10^{10}$ C3b: यानी
प्लाज़्मा के C3 का दस लाखवाँ भाग।
\textbf{12.} C3 के बिना तीनों पथ अपने साझा चरण पर रुक जाते हैं: न
\omterm{def:b3:innate-immunity:complement}{ऑप्सोनीकरण}, न \omterm{def:b3:innate-immunity:complement}{ऐनाफ़ाइलोटॉक्सिन}, न लयन --- संपुटित जीवाणुओं से गंभीर,
बार-बार लौटते संक्रमण। C9 के बिना केवल अंतिम छिद्र जाता है;
\omterm{def:b3:innate-immunity:complement}{ऑप्सोनीकरण} और \omterm{def:b3:innate-immunity:inflammation}{शोथ} काम करते हैं, और लक्षणरूप बार-बार लौटते
\emph{Neisseria} संक्रमण हैं, यानी वही एक वंश जिसे भक्षककोशिकाएँ ठीक
से नहीं सँभालतीं और लयन ठीक से सँभाल लेता है।
\textbf{13.} $2\times 0.24\times 8 = \qty{3.8}{MJ}$, लगभग \qty{100}{g}
वसा।
\textbf{14.} उस घंटे में एक द्विगुणन: $2\times 10^{4}$, न कि
$2.8\times 10^{4}$, यानी $1.4$ का गुणक।
\textbf{15.} घंटा 2: मारने से पहले $4\times 10^{4}$, बाद में एक भी नहीं
--- पहले की तरह \qty{2}{h} पर समाप्त; ज्वर का योगदान तब मायने रखता है
जब \omterm{def:b3:innate-immunity:innate}{न्यूट्रोफिलों} की पूर्ति सीमांत हो, जैसा प्रश्न 4 और 6 में।
\textbf{16.} $199\,\text{mg/L}\times 3\,\text{L} = \qty{0.6}{g}$; $0.6/
\num{115000} = 5.2\times 10^{-6}$ मोल $= 3.1\times 10^{18}$ अणु;
\qty{172800}{s} में, प्रति सेकंड $1.8\times 10^{13}$।
\textbf{17.} वे नियत बिंदु को वापस \qty{37}{\celsius} की ओर उतार देती
हैं, रोगी बेहतर महसूस करता है, और जीवाणु थोड़ा तेज़ दुगुने होते हैं; इस
प्रतिरूप में \omterm{def:b3:innate-immunity:innate}{न्यूट्रोफिल} फिर भी जीतते हैं, और यह प्रमाण कमज़ोर है कि
ज्वरनाशक साधारण संक्रमण बिगाड़ते हैं।
\textbf{18.} लगभग \qty{41}{\celsius} से ऊपर प्रोटीन विकृत होने लगते हैं,
एंज़ाइम विफल होती हैं और न्यूरॉन गड़बड़ चलते हैं (दौरे);
\qty{42}{\celsius} घातक है। नियत बिंदु उससे नीचे शीतजनकों ---
इंटरल्यूकिन-10, वैसोप्रेसिन, ग्लूकोकॉर्टिकॉइड --- से और अधश्चेतक में
प्रोस्टाग्लैंडिन-क्रिया की छत से थमा रहता है।
\textbf{19.} $10^{9}\times 3\times 10^{6} = 3\times 10^{15}$ अणु
$= 5\times 10^{-9}$ मोल, \qty{5}{L} में: \qty{1e-9}{mol/L}।
\textbf{20.} संतृप्तकारी सांद्रता के बराबर: शरीर का हर \omterm{def:b3:innate-immunity:innate}{महाभक्षक} एक साथ
अधिकतम सक्रिय है।
\textbf{21.} $10^{11}\times 1000\times 3600 = 3.6\times 10^{17}$ अणु
$= 6\times 10^{-7}$ मोल; \qty{15}{L} में, \qty{4e-8}{mol/L} --- यानी
सक्रिय सांद्रता से चार हज़ार गुना।
\textbf{22.} हर अणुधमनी फैलती है और हर अणुशिरा रिसती है: रक्त का आयतन
चौड़ी हुई वाहिकाओं में जमा हो जाता है और ऊतकों में निकल भागता है,
इसलिए दबाव गिरता है; गुर्दे, कम सिंचित, छानना बंद कर देते हैं; और
हज़ारों केशिकाओं में छिड़ा स्कंदन उन्हें अवरुद्ध कर देता है तथा स्कंदन
कारक चुका देता है।
\textbf{23.} मारे गए \omterm{def:b3:bacteriology:envelope}{ग्राम-ऋणात्मक} अपना सारा \omterm{def:b3:bacteriology:envelope}{लिपोपॉलिसैकेराइड} एक साथ
छोड़ देते हैं, यानी TLR4 के लिए कोई बड़ी खुराक, इसलिए \omterm{def:b3:innate-immunity:inflammation}{साइटोकाइन} का उछाल
घंटों तक बिगड़ सकता है। प्रति-TNF इसलिए विफल होता है कि आघात के समय तक
वह शृंखला TNF से आगे जा चुकी होती है --- इंटरल्यूकिन-1,
इंटरल्यूकिन-6 और स्कंदन पहले ही चल रहे होते हैं --- और इसलिए भी कि
जीवाणुओं पर काबू रखने के लिए TNF चाहिए: समस्या लक्ष्य की नहीं, समय की
है।
\textbf{24.} अंतःक्षिप्त \omterm{def:b3:bacteriology:envelope}{लिपोपॉलिसैकेराइड} TLR4-रहित चूहे का कुछ नहीं
बिगाड़ता, जो सामान्य चूहों के लिए घातक मात्राएँ झेल जाता है; कोई जीवित
\omterm{def:b3:bacteriology:envelope}{ग्राम-ऋणात्मक} संक्रमण उसे मार देता है, क्योंकि वह जीवाणुओं को जल्दी
पकड़ नहीं सकता और वह \omterm{def:b3:innate-immunity:inflammation}{शोथ} खड़ा नहीं कर सकता जो उन्हें घेरता है। यह
विरोधाभास घुल जाता है: जो संसूचक आघात लाता है वही संसूचक बचाव भी देता
है।
\textbf{25.} \qty{6}{h} पर कोई जीवाणु नहीं बचता (सारे \qty{2}{h} तक मारे
गए); \qty{120}{s} पर प्रति जीवाणु लगभग $1.8\times 10^{5}$ C3b; और दो
दिन का ज्वर \qty{3.8}{MJ} लेता है, यानी कोई \qty{100}{g} वसा।
\end{solution}
